% GENERATED FILE. Edit canonical lessons, then run npm run generate:books.
% canonical-source-hash: fnv1a64:8d0f3deedc85abfc
% canonical-lessons: ES-C666-saco-de-dormir, ES-C666-acampar, ES-C666-esterilla, ES-C666-caravana, ES-C666-linterna

\chapter{Bajo las estrellas}
\label{ch:es-bajo-las-estrellas}
\hlchaptermodality{666}

\begin{chapteropening}
\textbf{By the end of this chapter:} \emph{I can go camping: a sleeping bag and a mat, travelling in a caravan and bringing a torch.}

Complete the last lesson of chapter 666: I can go camping: a sleeping bag and a mat, travelling in a caravan and bringing a torch.
\end{chapteropening}

\section[el saco de dormir]{el saco de dormir — the sleeping bag}

\label{lesson:ES-C666-saco-de-dormir}



\begin{quote}
Before the new one: say the Spanish for the manager, then the Spanish for the owner.
\end{quote}

\subsection*{You'll want to know: el saco de dormir}
\textbf{el saco de dormir} — \textquotedblleft{}the sleeping bag\textquotedblright{}.

\begin{quote}
Hace frío, lleva un buen saco de dormir.
\end{quote}

\begin{quote}
\emph{It is cold, take a good sleeping bag.}
\end{quote}

In Latin America also \emph{la bolsa de dormir}.

\begin{cousinweb}
\emph{saco} is inherited from Latin \emph{saccus} \textquotedblleft{}a sack, a bag\textquotedblright{} (Lewis \& Short; EtymDB), and \emph{dormir} from Latin \emph{dormīre} \textquotedblleft{}to sleep\textquotedblright{}.
\end{cousinweb}

\begin{grammarlens}[title={What you've built}]
One more word for telling, explaining and describing at B1.
\end{grammarlens}

\subsection*{Your turn}
\begin{itemize}
\raggedright
  \item \emph{Say it:} \emph{el saco de dormir}
  \item \emph{Say it:} \emph{el saco de dormir}, once more
  \item \emph{Say it:} a sentence of your own with \emph{el saco de dormir}
  \item \emph{Recall:} say the Spanish for the manager, then the Spanish for the owner, then say \emph{el saco de dormir} again
\end{itemize}

\begin{tcolorbox}[breakable,colback=teal!4,colframe=teal!35!black,title={Before you move on}]
What does el saco de dormir mean? (\textquotedblleft{}The sleeping bag\textquotedblright{}.) Say it once more.
\end{tcolorbox}
\section[acampar]{acampar — to camp}

\label{lesson:ES-C666-acampar}



\begin{quote}
Before the new one: say the Spanish for the owner, then the Spanish for the sleeping bag.
\end{quote}

\subsection*{You'll want to know: acampar}
\textbf{acampar} — \textquotedblleft{}to camp\textquotedblright{}.

\begin{quote}
Está prohibido acampar en la playa.
\end{quote}

\begin{quote}
\emph{Camping on the beach is forbidden.}
\end{quote}

\emph{La acampada libre} — wild camping.

\begin{cousinweb}
\emph{acampar} was borrowed from Italian \emph{accampare} (EtymDB), built on \emph{campo} \textquotedblleft{}field\textquotedblright{}, from Latin \emph{campus}.
\end{cousinweb}

\begin{grammarlens}[title={What you've built}]
One more word for telling, explaining and describing at B1.
\end{grammarlens}

\subsection*{Your turn}
\begin{itemize}
\raggedright
  \item \emph{Say it:} \emph{acampar}
  \item \emph{Say it:} \emph{acampar}, once more
  \item \emph{Say it:} a sentence of your own with \emph{acampar}
  \item \emph{Recall:} say the Spanish for the owner, then the Spanish for the sleeping bag, then say \emph{acampar} again
\end{itemize}

\begin{tcolorbox}[breakable,colback=teal!4,colframe=teal!35!black,title={Before you move on}]
What does acampar mean? (\textquotedblleft{}To camp\textquotedblright{}.) Say it once more.
\end{tcolorbox}
\section[la esterilla]{la esterilla — the sleeping mat; the small mat}

\label{lesson:ES-C666-esterilla}



\begin{quote}
Before the new one: say the Spanish for the sleeping bag, then the Spanish for to camp.
\end{quote}

\subsection*{You'll want to know: la esterilla}
\textbf{la esterilla} — \textquotedblleft{}the sleeping mat; the small mat\textquotedblright{}.

\begin{quote}
Pongo la esterilla debajo del saco.
\end{quote}

\begin{quote}
\emph{I put the mat under the sleeping bag.}
\end{quote}

At the gym or the beach, \emph{una esterilla} is a mat.

\begin{cousinweb}
\emph{esterilla} is a little \emph{estera} \textquotedblleft{}mat\textquotedblright{}, which is inherited from Latin \emph{storea} \textquotedblleft{}a mat of plaited straw or rushes\textquotedblright{} (Lewis \& Short; EtymDB).
\end{cousinweb}

\begin{grammarlens}[title={What you've built}]
One more word for telling, explaining and describing at B1.
\end{grammarlens}

\subsection*{Your turn}
\begin{itemize}
\raggedright
  \item \emph{Say it:} \emph{la esterilla}
  \item \emph{Say it:} \emph{la esterilla}, once more
  \item \emph{Say it:} a sentence of your own with \emph{la esterilla}
  \item \emph{Recall:} say the Spanish for the sleeping bag, then the Spanish for to camp, then say \emph{la esterilla} again
\end{itemize}

\begin{tcolorbox}[breakable,colback=teal!4,colframe=teal!35!black,title={Before you move on}]
What does la esterilla mean? (\textquotedblleft{}The sleeping mat; the small mat\textquotedblright{}.) Say it once more.
\end{tcolorbox}
\section[la caravana]{la caravana — the caravan, the camper; the tailback}

\label{lesson:ES-C666-caravana}



\begin{quote}
Before the new one: say the Spanish for to camp, then the Spanish for the sleeping mat.
\end{quote}

\subsection*{You'll want to know: la caravana}
\textbf{la caravana} — \textquotedblleft{}the caravan, the camper; the tailback\textquotedblright{}.

\begin{quote}
Viajamos por la costa en caravana.
\end{quote}

\begin{quote}
\emph{We travelled along the coast in a caravan.}
\end{quote}

In Spain \emph{una caravana} is also a long line of cars.

\begin{cousinweb}
\emph{caravana} comes from Persian \emph{kārvān} \textquotedblleft{}a company of travellers\textquotedblright{} (EtymDB), which reached Spanish through French or Italian.
\end{cousinweb}

\begin{grammarlens}[title={What you've built}]
One more word for telling, explaining and describing at B1.
\end{grammarlens}

\subsection*{Your turn}
\begin{itemize}
\raggedright
  \item \emph{Say it:} \emph{la caravana}
  \item \emph{Say it:} \emph{la caravana}, once more
  \item \emph{Say it:} a sentence of your own with \emph{la caravana}
  \item \emph{Recall:} say the Spanish for to camp, then the Spanish for the sleeping mat, then say \emph{la caravana} again
\end{itemize}

\begin{tcolorbox}[breakable,colback=teal!4,colframe=teal!35!black,title={Before you move on}]
What does la caravana mean? (\textquotedblleft{}The caravan, the camper; the tailback\textquotedblright{}.) Say it once more.
\end{tcolorbox}
\section[la linterna]{la linterna — the torch, the flashlight}

\label{lesson:ES-C666-linterna}



\begin{quote}
Before the new one: say the Spanish for the sleeping mat, then the Spanish for the caravan.
\end{quote}

\subsection*{You'll want to know: la linterna}
\textbf{la linterna} — \textquotedblleft{}the torch, the flashlight\textquotedblright{}.

\begin{quote}
Lleva una linterna para la noche.
\end{quote}

\begin{quote}
\emph{Take a torch for the night.}
\end{quote}

A lantern on a ship or a tower is also \emph{la linterna}.

\begin{cousinweb}
\emph{linterna} comes from older \emph{lanterna} (EtymDB), a learned borrowing from Latin \emph{lanterna} \textquotedblleft{}a lantern, a lamp, a torch\textquotedblright{} (Lewis \& Short).
\end{cousinweb}

\begin{grammarlens}[title={What you've built}]
One more word for telling, explaining and describing at B1.
\end{grammarlens}

\subsection*{Your turn}
\begin{itemize}
\raggedright
  \item \emph{Say it:} \emph{la linterna}
  \item \emph{Say it:} \emph{la linterna}, once more
  \item \emph{Say it:} a sentence of your own with \emph{la linterna}
  \item \emph{Recall:} say the Spanish for the sleeping mat, then the Spanish for the caravan, then say \emph{la linterna} again
\end{itemize}

\begin{tcolorbox}[breakable,colback=teal!4,colframe=teal!35!black,title={Before you move on}]
What does la linterna mean? (\textquotedblleft{}The torch, the flashlight\textquotedblright{}.) Say it once more.
\end{tcolorbox}
